% =============================================================================
\chapter{Riscos e Estratégias de Mitigação}
% =============================================================================

\section{Escopo excessivo}

Integrar HTN, geração de preferências, MORL, simulação rica e coordenação
multiagente pode exceder o escopo do mestrado. A contribuição principal será
restrita à seleção online de métodos HTN para um agente. Multiagente e gerador
de pesos aprendido permanecem opcionais.

\section{Recompensas esparsas ou atrasadas}

Uma escolha hierárquica pode afetar resultados somente após várias tarefas
primitivas. Serão avaliadas recompensas intermediárias por objetivo, janelas de
acumulação e estratégias de \textit{reward shaping}, sempre documentando o
efeito sobre a política.

\section{Aprendizado instável sob preferências variadas}

A política pode não generalizar em todo o simplex de pesos. A mitigação inclui
começar com poucos objetivos, normalizar escalas, usar amostragem controlada e
comparar desempenho em vetores vistos e não vistos.

\section{Oscilação comportamental}

Mudanças frequentes de preferência podem causar trocas repetidas de método.
Serão considerados períodos mínimos de compromisso, custo de troca, histerese
ou limiares de decisão. Essas intervenções serão tratadas como parte explícita
do desenho experimental.

\section{Ambiguidade na atribuição de crédito}

A duração variável das decomposições pode dificultar associar retornos à
escolha do método. O protocolo definirá claramente início, término,
interrupção e desconto da decisão hierárquica, com testes em cenários curtos
antes do domínio completo.

\section{Novidade científica insuficiente}

Há trabalhos que combinam planejamento e aprendizado por reforço. A mitigação
é realizar revisão estruturada e posicionar a contribuição de forma precisa na
interseção entre filtragem HTN, preferências dinâmicas e MORL condicionado para
seleção online de métodos.

\section{Comparação experimental injusta}

Diferenças de orçamento, ajuste ou dificuldade podem favorecer uma
arquitetura. Serão usados cenários pareados, mesmas sementes de avaliação,
orçamentos comparáveis, critérios explícitos de ajuste e relato de tamanhos de
efeito.
